\documentclass{beamer}

\usepackage{../tema}


\title{SAL}
\author[F. Pasqual, C. Libralato]{F. Pasqual, C. Libralato}
\setDate{11/03/2026}

\begin{document}

{
\setbeamertemplate{footline}{}
\begin{frame}
  \titlepage
\end{frame}
}

\addMember{Filippo}{Venzo}{Amministratore}
\addMember{Giuseppe}{De Fina}{Progettista}
\addMember{Valeria}{Baleanu}{Progettista}
\addMember{Christian}{Libralato}{Verificatore}
\addMember{Leonardo}{Pellizzon}{Progettista}
\addMember{Francesco}{Pasqual}{Responsabile}
\addMember{Luca}{Granziero}{Progettista}
\makeMembersFrame

\begin{frame}{Avanzamenti settimanali}
  \begin{itemize}
    \item Continuazione della progettazione, coinvolgimento di tutto il team;
    \item Definizione dei primi bozzetti \textit{wireframe$_G$} su Figma.
  \end{itemize}
\end{frame}

% PRIMO FRAME: Attività terminate dal verbale esterno 
\addPrecedentTodo{ve\_2026\_03\_04.a2}{-}{Proseguire le attività di progettazione tenendo conto dei chiarimenti emersi durante la riunione e delle nuove informazioni acquisite}{F. Venzo, V. Baleanu, F. Pasqual, G. De Fina}{08/03/2026}


% Reset delle entries per il secondo frame 
\renewcommand{\precedentTodoEntries}{}

%  SECONDO FRAME: Attività terminate dal verbale interno
\addPrecedentTodo{vi\_2026\_02\_23.a2 }{\#143}{Aggiornamento PdP a sprint 8}{L. Pellizzon}{08/03/2026}

\makePrecedentTodoFrame


\begin{frame}{Attività da completare}
  \begin{itemize}
    \item Aggiornare Piano di Progetto e Piano di Qualifica allo sprint 9;
    \item Disporre di una progettazione a un livello di dettaglio adeguato per poter avviare la programmazione durante questo sprint.
  \end{itemize}
\end{frame}

\begin{frame}{Dubbi incontrati}
  \begin{itemize}
    \item Separazione dei flussi di autenticazione e token durante la progettazione del \textit{frontend};
    \item Caratteristiche, numerosità, e livello di dettaglio dei bozzetti \textit{wireframe}
  \end{itemize}
\end{frame}

\end{document}